\textbf{Phần 1. Trắc nghiệm lựa chọn}
\setcounterpageref{ex}{0}

\begin{ex} %[5D6N30]
Số thập phân gồm \textit{Hai mươi ba đơn vị, tám mươi mốt phần trăm} được viết là
\choice
{$238{,}1$}
{$2{,}381$}
{\True $23{,}81$}
{$0{,}2381$}
\loigiai{
Số thập phân đó là $23{,}81$.\\
\textbf{Đáp án của câu hỏi là C}
}
\end{ex}

\begin{ex} %[5D6N30]
Chữ số $3$ trong số thập phân $2{,}03$ có giá trị thuộc hàng nào?
\choice
{\True Hàng phần trăm}
{Hàng chục}
{Hàng đơn vị}
{Hàng phần mười}
\loigiai{
Chữ số $3$ đứng ở vị trí thứ hai sau dấu phẩy nên thuộc hàng phần trăm.\\
\textbf{Đáp án của câu hỏi là A}
}
\end{ex}

\begin{ex} %[5D6H30]
Đổi $8\text{ mm}$ ra mét ta được số thập phân là
\choice
{$0{,}8\text{ m}$}
{$0{,}08\text{ m}$}
{$8\,000\text{ m}$}
{\True $0{,}008\text{ m}$}
\loigiai{
Vì $1\text{ m} = 1\,000\text{ mm}$ nên $8\text{ mm} = \dfrac{8}{1\,000}\text{ m} = 0{,}008\text{ m}$.\\
\textbf{Đáp án của câu hỏi là D}
}
\end{ex}

\begin{ex} %[5D6H30]
So sánh hai số thập phân $17{,}1$ và $9{,}725$, khẳng định nào sau đây là đúng?
\choice
{$17{,}1 < 9{,}725$ vì $17{,}1$ có $3$ chữ số còn $9{,}725$ có $4$ chữ số}
{\True $17{,}1 > 9{,}725$ vì phần nguyên $17 > 9$}
{$17{,}1 = 9{,}725$}
{Không so sánh được vì số chữ số khác nhau}
\loigiai{
So sánh phần nguyên trước: $17 > 9$ nên $17{,}1 > 9{,}725$.\\
\textbf{Đáp án của câu hỏi là B}
}
\end{ex}

\begin{ex} %[5D6H30]
Đổi $800\text{ kg}$ ra tấn ta được số thập phân là
\choice
{$8\text{ tấn}$}
{$0{,}08\text{ tấn}$}
{\True $0{,}8\text{ tấn}$}
{$80\text{ tấn}$}
\loigiai{
Vì $1\text{ tấn} = 1\,000\text{ kg}$ nên $800\text{ kg} = \dfrac{800}{1\,000}\text{ tấn} = 0{,}8\text{ tấn}$.\\
\textbf{Đáp án của câu hỏi là C}
}
\end{ex}

\begin{ex} %[5D6H30]
Viết số đo diện tích $1\text{ m}^2\,5\text{ dm}^2$ dưới dạng số thập phân với đơn vị mét vuông là
\choice
{\True $1{,}05\text{ m}^2$}
{$1{,}5\text{ m}^2$}
{$1{,}005\text{ m}^2$}
{$15\text{ m}^2$}
\loigiai{
$1\text{ m}^2\,5\text{ dm}^2 = 1\dfrac{5}{100}\text{ m}^2 = 1{,}05\text{ m}^2$.\\
\textbf{Đáp án của câu hỏi là A}
}
\end{ex}

\begin{ex} %[5D6H30]
Làm tròn số thập phân $9{,}548$ đến hàng phần trăm ta được số
\choice
{$9{,}5$}
{$9{,}54$}
{\True $9{,}55$}
{$10$}
\loigiai{
Chữ số ở hàng phần nghìn là $8 \ge 5$ nên ta tăng chữ số hàng phần trăm thêm $1$: được $9{,}55$.\\
\textbf{Đáp án của câu hỏi là C}
}
\end{ex}

\begin{ex} %[5D6H30]
Kỉ lục chạy $100\text{ m}$ là $9{,}58\text{ giây}$. Làm tròn kết quả này đến hàng phần mười ta được
\choice
{$9{,}5\text{ giây}$}
{$10\text{ giây}$}
{$9{,}58\text{ giây}$}
{\True $9{,}6\text{ giây}$}
\loigiai{
Chữ số hàng phần trăm là $8 \ge 5$ nên làm tròn lên được $9{,}6\text{ giây}$.\\
\textbf{Đáp án của câu hỏi là D}
}
\end{ex}

\begin{ex} %[5D6V30]
Cân nặng của $4$ loại quả trong siêu thị là: Táo nặng $2{,}5\text{ kg}$, dưa hấu nặng $4{,}19\text{ kg}$, chuối nặng $4{,}0\text{ kg}$, nho nặng $1{,}75\text{ kg}$. Loại quả nhẹ nhất là
\choice
{Táo}
{\True Nho}
{Chuối}
{Dưa hấu}
\loigiai{
So sánh: $1{,}75\text{ kg} < 2{,}5\text{ kg} < 4{,}0\text{ kg} < 4{,}19\text{ kg}$. Vậy quả nhẹ nhất là nho.\\
\textbf{Đáp án của câu hỏi là B}
}
\end{ex}

\begin{ex} %[5D6V30]
Tìm chữ số thích hợp thay vào dấu $?$ để được số $2{,}?1$ thỏa mãn $2 < 2{,}?1 < 3$.
\choice
{Chỉ có thể là chữ số $0$}
{Chỉ có thể là chữ số $9$}
{Không có chữ số nào thỏa mãn}
{\True Bất kì chữ số nào từ $0$ đến $9$}
\loigiai{
Với mọi chữ số $?$ từ $0$ đến $9$, phần nguyên đều bằng $2$ và phần thập phân $> 0$ nên $2 < 2{,}?1 < 3$ luôn đúng.\\
\textbf{Đáp án của câu hỏi là D}
}
\end{ex}

\begin{ex} %[5D6V30]
Một túi cà chua cân nặng $1{,}45\text{ kg}$ và một túi hành tây cân nặng $500\text{ g}$. Cả hai túi cân nặng tất cả là
\choice
{$1{,}5\text{ kg}$}
{\True $1{,}95\text{ kg}$}
{$6{,}45\text{ kg}$}
{$1{,}455\text{ kg}$}
\loigiai{
Đổi $500\text{ g} = 0{,}5\text{ kg}$. Tổng cân nặng là: $1{,}45 + 0{,}5 = 1{,}95\text{ (kg)}$.\\
\textbf{Đáp án của câu hỏi là B}
}
\end{ex}

\begin{ex} %[5D6V30]
Sắp xếp các số đo chiều cao: Mai cao $150\text{ cm}$, Mi cao $1{,}2\text{ m}$, Rô-bốt cao $90\text{ cm}$. Thứ tự các bạn từ cao nhất đến thấp nhất là
\choice
{\True Mai, Mi, Rô-bốt}
{Mi, Mai, Rô-bốt}
{Rô-bốt, Mi, Mai}
{Mai, Rô-bốt, Mi}
\loigiai{
Đổi về cùng đơn vị mét: Mai cao $1{,}5\text{ m}$; Mi cao $1{,}2\text{ m}$; Rô-bốt cao $0{,}9\text{ m}$.\\
Thứ tự từ cao đến thấp: Mai ($1{,}5\text{ m}$) $>$ Mi ($1{,}2\text{ m}$) $>$ Rô-bốt ($0{,}9\text{ m}$).\\
\textbf{Đáp án của câu hỏi là A}
}
\end{ex}

\textbf{Phần 2. Trắc nghiệm đúng sai}
\setcounterpageref{ex}{0}

\begin{ex} %[5D6H30]
Xét tính đúng sai của các khẳng định sau về số thập phân:
\choiceTF
{\True Số $100{,}10$ bằng số $100{,}1$.}
{Chữ số $3$ trong số $109{,}37$ thuộc hàng chục.}
{\True Đổi $17\text{ ml} = 0{,}017\text{ l}$.}
{\True Làm tròn số $100{,}917$ đến hàng phần trăm ta được $100{,}92$.}
\loigiai{
\begin{itemize}
    \item $100{,}10 = 100{,}1$ (bỏ chữ số 0 tận cùng bên phải phần thập phân).
    \item Chữ số 3 trong $109{,}37$ thuộc hàng phần mười, không phải hàng chục.
    \item $17\text{ ml} = 0{,}017\text{ l}$.
    \item $100{,}917$ có chữ số hàng phần nghìn là $7 \ge 5$ nên làm tròn thành $100{,}92$.
\end{itemize}
\begin{itemize}
    \item a) \textbf{Đúng}.
    \item b) \textbf{Sai}.
    \item c) \textbf{Đúng}.
    \item d) \textbf{Đúng}.
\end{itemize}
}
\end{ex}

\begin{ex} %[5D6V30]
Kỉ lục thế giới chạy $200\text{ m}$ là $19{,}19\text{ giây}$. Xét tính đúng sai của các khẳng định sau:
\choiceTF
{\True Phần nguyên của số chỉ thời gian kỉ lục là $19$.}
{\True Làm tròn kỉ lục đến hàng phần mười ta được $19{,}2\text{ giây}$.}
{Làm tròn kỉ lục đến số tự nhiên gần nhất ta được $20\text{ giây}$.}
{\True Thời gian kỉ lục trên ít hơn $20\text{ giây}$.}
\loigiai{
\begin{itemize}
    \item Phần nguyên là 19.
    \item Chữ số hàng phần trăm là $9 \ge 5 \Rightarrow$ làm tròn đến hàng phần mười là $19{,}2\text{ s}$.
    \item Chữ số hàng phần mười là $1 < 5 \Rightarrow$ làm tròn đến số tự nhiên gần nhất là $19\text{ giây} \neq 20\text{ giây}$.
    \item $19{,}19 < 20$.
\end{itemize}
\begin{itemize}
    \item a) \textbf{Đúng}.
    \item b) \textbf{Đúng}.
    \item c) \textbf{Sai}: Làm tròn là $19\text{ giây}$.
    \item d) \textbf{Đúng}.
\end{itemize}
}
\end{ex}

\textbf{Phần 3. Trả lời ngắn}
\setcounterpageref{ex}{0}

\begin{ex} %[5D6H30]
Đổi $173\text{ cm}$ ra mét dưới dạng số thập phân.
\shortans{1,73}
\loigiai{
\textbf{Đáp số:} $1{,}73$\\
$173\text{ cm} = 1{,}73\text{ m}$.
}
\end{ex}

\begin{ex} %[5D6H30]
Đổi $3\text{ dm}^2\,75\text{ cm}^2$ ra đề-xi-mét vuông dưới dạng số thập phân.
\shortans{3,75}
\loigiai{
\textbf{Đáp số:} $3{,}75$\\
$3\text{ dm}^2\,75\text{ cm}^2 = 3{,}75\text{ dm}^2$.
}
\end{ex}

\begin{ex} %[5D6H30]
Làm tròn số thập phân $17{,}153$ đến hàng phần trăm. (Nhập kết quả dưới dạng số thập phân).
\shortans{17,15}
\loigiai{
\textbf{Đáp số:} $17{,}15$\\
Chữ số hàng phần nghìn là $3 < 5$ nên giữ nguyên: $17{,}15$.
}
\end{ex}

\begin{ex} %[5D6V30]
Một túi cà chua nặng $1{,}45\text{ kg}$ và túi hành tây nặng $0{,}5\text{ kg}$. Tính tổng cân nặng của hai túi theo ki-lô-gam. (Nhập kết quả dưới dạng số thập phân).
\shortans{1,95}
\loigiai{
\textbf{Đáp số:} $1{,}95$\\
$1{,}45 + 0{,}5 = 1{,}95\text{ (kg)}$.
}
\end{ex}

\begin{ex} %[5D6V30]
Kỉ lục chạy $100\text{ m}$ là $9{,}58\text{ giây}$. Làm tròn số đo thời gian này đến hàng phần mười. (Nhập kết quả dưới dạng số thập phân).
\shortans{9,6}
\loigiai{
\textbf{Đáp số:} $9{,}6$\\
$9{,}58 \approx 9{,}6\text{ (giây)}$.
}
\end{ex}

\begin{ex} %[5D6V30]
Đổi $3\text{ dm}^2$ ra mét vuông dưới dạng số thập phân.
\shortans{0,03}
\loigiai{
\textbf{Đáp số:} $0{,}03$\\
$3\text{ dm}^2 = 0{,}03\text{ m}^2$.
}
\end{ex}
